% 3 Pages
\section{Experimental Results}
\label{sec_experiment}

% --------------------------------------------------------
% \subsection{Experiment settings}
% \label{sec_Implementation_Details}

% --------------------------------------------------------
\noindent \textbf{Datasets.}
To probe the effectiveness of our proposed framework, {\papernameAbbrev}, we conduct evaluations on three 3D visual-grounding benchmarks, namely Nr3D, Sr3D \citep{achlioptas2020referit3d} and ScanRefer \citep{chen2020scanrefer}.
Nr3D contains 41.5K natural, free-form utterances gathered from humans through a referring game, while Sr3D consists of 83.5K synthetic utterances. 
Consequently, ScanRefer provides 51.5K utterances of 11K objects for 800 3D indoor scenes.
% --------------------------------------------------------

\noindent \textbf{Network Configuration.}
%To prove the compatibility of our proposed framework, {\papernameAbbrev}, we integrated it into two different architectures; MVT \citep{huang2022multi} and ViL3DRel \citep{chen2022language}.
%We integrate our CoT module into the SOTA architecture MVT \citep{huang2022multi}.
We model the Pathway module using only one transformer encoder layer, and the CoT decoder using a single transformer decoder layer.
The number of heads used are 7 and 16 for the Pathway module and CoT decoder, respectively.
The number of proposals $L$ and the maximum sentence length $W$ are 52 and 24, respectively. $L$ and $W$ define the sizes of the input sequences to our CoT decoder.
The maximum number $M$ of objects in the sentence, the output sequence length for our CoT decoder, is 8 and 3 for Nr3D and Sr3D, respectively.
Following previous works \citep{abdelreheem20223dreftransformer, achlioptas2020referit3d, he2021transrefer3d, roh2022languagerefer, jain2021looking, yang2021sat, qi2017pointnet++}, we randomly sample 1024 points for each proposal, set the hidden dimensions $d$ to 768, and train the model for 100 epochs from scratch using the weight initialization strategy described in \citep{he2015delving}.
The initial learning rate is set to $10^{-4}$ and decreases by 0.65 every ten epochs. 
The Adam optimizer \citep{kingma2014adam} and a mini-batch size of 24 per GPU are used for training all the models.
We set the losses weights as follows: $\lambda_{{\mathcal{V}}}=5$, $\lambda_{{\mathcal{T}}}=0.5$, $\mathcal{L}_{ref}=5$, and $\lambda_{dist}=1$.
We used the PyTorch framework and a single NVIDIA A6000 GPU for training.
% --------------------------------------------------------
\subsection{Ablation Studies}

We conducted several ablation studies to validate each module in our framework, termed {\papernameAbbrev}.

% -------------------------------------------------------
\begin{table}[!tb]
\vspace{-8mm}
\tabcolsep=0.02cm
\begin{minipage}[c]{0.55\textwidth}
\captionsetup{font=scriptsize}
\centering
\captionof{table}{Ablation study for different components of our {\papernameAbbrev} framework. First, we compare the baseline, i.e., MVT \citep{huang2022multi}, against the parallel and the chain-of-thoughts approaches. Then, we show the effect of adding distractor loss. All the experiments are conducted using the Nr3D dataset \citep{achlioptas2020referit3d}.}
\scalebox{0.8}{
\begin{tabular}{cccccccc}
    \toprule
    \- &  \makecell{Data\\Percentage} & \makecell{+Distractor.\\ Loss.} & \makecell{+Parallel} & \makecell{+CoT} & \- & Nr3D ↑ & Sr3D ↑ \\
    \cmidrule(r){1-5}
    \cmidrule(r){7-8}
     (a) &  10\%      &  -      &  -      & -       & \- & 27.6 & 48.8 \\ 
     (b) &  10\%      &  -      & \cmark  & -       & \- & 31.7 & 55.3 \\ 
     (c) &  10\%      &  -      & -       & \cmark  & \- & 37.5 & 65.2 \\ 
     (d) &  10\%      &  \cmark & -       & \cmark  & \- & \textbf{38.2} & \textbf{66.4} \\ 
     \midrule
     (e) &  100\%     &  -      &  -      & -       & \- & 55.2 & 66.0 \\ 
     (f) &  100\%     &  -      & \cmark  & -       & \- & 57.0 & 71.5 \\ 
     (g) &  100\%     &  -      & -       & \cmark  & \- & 60.0 & 72.7 \\ 
     (h) &  100\%     &  \cmark & -       & \cmark  & \- & \textbf{60.4} & \textbf{73.2} \\ 
    \bottomrule
  \end{tabular}
}
\label{Table_ablation_for_cot_components}
\end{minipage}
\hfil
\hfil
\hfil
\begin{minipage}[c]{0.42\textwidth}
    \centering
    \captionsetup{font=scriptsize}
    \captionof{table}{Ablation study for different numbers of the transformer blocks used in our CoT decoder, based on MVT.}
    \label{Table_ablation_transformer_blocks}
    \scalebox{0.65}{
    \begin{tabular}{cccc}
        \toprule[1.5pt]
        \textbf{\makecell{\# Transformer \\Blocks}} & \- & \textbf{Nr3D ↑} & \textbf{Sr3D ↑} \\
        \cmidrule(r){1-1}
        \cmidrule(r){3-4}
        \textbf{1} & \- & \textbf{60.4} & \textbf{73.2} \\ 
        \textbf{2} & \- & 60.4 & 73.3 \\ 
        \textbf{4} & \- & 60.1 & 72.9 \\ 
        \bottomrule[1.5pt]
    \end{tabular}
    }
    % Add some space between the two right tables
    %\vspace{2pt}

    % Lower right table
    \captionsetup{font=scriptsize}
    \captionof{table}{Ablation study to highlight the effect of the anchors' quality on the final target referring accuracy, based on Nr3D dataset and MVT as a baseline.}
    \label{Table_ablation_anchors_quality}
    \scalebox{0.60}{
    \begin{tabular}{cccc}
        \toprule[1.5pt]
        \multirow{2}{*}{\textbf{Method}} & \- & \multirow{2}{*}{\makecell{\textbf{Anchors}\\\textbf{Ref. Acc.}}} & \multirow{2}{*}{\makecell{\textbf{Target}\\\textbf{Ref. Acc.}}} \\
        \- & \-&\-& \- \\
        \cmidrule(r){1-1}
        \cmidrule(r){3-4}
        Baseline & \- & N/A & 55.1 \\
        +CoT +Zeroing Anchor Loss & \- & 4.5 & 55.0 \\
        +CoT +Pseudo Labels & \- & 72.3 & 60.4 \\
        +CoT +Human Labels & \- & \textbf{73.6} & \textbf{64.4} \\
        \bottomrule[1.5pt]
    \end{tabular}
    }
\end{minipage}
\end{table}
% -------------------------------------------------------

\noindent \textbf{CoT vs. Parallel.}
To assess our {\papernameAbbrev} framework, we have to disentangle the CoT from the pseudo label generation module.
In other words, it could be thought that the achieved gain caused by just accessing additional supervision signal, i.e., anchors' annotations.
To this end, we implement a parallel approach, that has access to the anchors' labels, however, localizes the targets and anchors in one shot without any interaction between them.
In contrast, our {\papernameAbbrev} framework leverages the causality between the anchors and the target through our chain-of-thoughts decoder.
As shown in Table \ref{Table_ablation_for_cot_components}, on the challenging setup, where we assume access for only 10\% of the training data while testing on the entire testing dataset, the parallel variant boosts the performance by 4\% and 6.5\% over the vanilla MVT using Nr3D and Sr3D, respectively (row b).
On the other hand, our {\papernameAbbrev} framework surpasses the vanilla MVT by 10\% and 16.4\% using Nr3D and Sr3D, respectively (row c).
Consequently, using the entire data, our CoT surpasses both the parallel and the vanilla MVT approaches by 3\% and 5\% on Nr3D, and by 1\% and 6.7\% on Sr3D, respectively.

\noindent \textbf{Distractor Loss.}
We add an auxiliary distractor binary classification loss $\mathcal{L}_{dist}$ to disambiguate the target and the distractors, as discussed in Sec. \ref{sec_loss}.
As shown in Table \ref{Table_ablation_for_cot_components}, rows d and h, incorporating it boosts the performance by 0.5-1\%.

\noindent \textbf{Anchors Quality Effect.}
To establish the affirmative role of anchors in refining target localization accuracy without inducing detrimental effects, we conducted a worst-case simulation. In this simulation, all anchors were falsely detected. 
In other words, we deliberately zero the localization loss associated with anchors during training, causing the decoder to randomly and inaccurately predict anchor locations.
Remarkably, the target accuracy remained unaltered, as shown in Table \ref{Table_ablation_anchors_quality}, reflecting the robustness of the approach. 
The accuracy experienced a decrement of approximately 5\%, dropping from 60.4\% to 55\%. Encouragingly, even in this scenario, the accuracy mirrored the baseline performance, steadfast at 55.1\%. This substantiates that while anchor detection may exhibit inaccuracies, the broader framework's efficacy in target localization remains largely unaffected.
In contrast, we replaced the pseudo labels with manual annotations \citep{abdelreheem2022scanents3d}. 
This substitution serves as an upper-bound reference point for evaluation. 
As shown in Table \ref{Table_ablation_anchors_quality}, the exchange of noisy pseudo labels with precise manual annotations led to a noteworthy 4\% enhancement in referring accuracy for Nr3D, elevating it from 60.4\% to 64.4\%.

\noindent \textbf{Number of Transformer Blocks.}
As shown in Table \ref{Table_ablation_transformer_blocks}, we have explored using 1, 2, and 4 transformer blocks in our CoT referring decoder. 
However, we didn't notice a significant gain in performance; therefore, we preferred to use a single transformer block.
% --------------------------------------------------------
\subsection{Comparison to State-of-the-art}

% -------------------------------------------------------
\tabcolsep=0.1cm
\begin{table}
\vspace{-8mm}
\captionsetup{font=small}
\centering
\caption{Benchmarking results on Nr3D and Sr3D datasets \citep{achlioptas2020referit3d}. We emphasize that we did not require any additional GT annotations, in contrast to SAT \citep{yang2021sat} which requires access to real 2D images, and PhraseRefer \citep{yuan2022toward} and ScanEnts \citep{abdelreheem2022scanents3d} require manual annotations for the whole anchors in the data. We report the standard-deviation $\sigma$ in green.}
\scalebox{0.6}{
\begin{tabular}{lcccccccccccccc}
    \toprule
    \multicolumn{1}{c}{\multirow{2}{*}{Method}} & \- & \multicolumn{1}{c}{\multirow{2}{*}{\makecell{GT \\ Anchors}}} & \- &  \multicolumn{5}{c}{Nr3D ↑} & \- &\multicolumn{5}{c}{Sr3D ↑}\\
    %\cmidrule(r){1-2}
    \cmidrule(r){5-9}
    \cmidrule(r){10-15}
    \- & \- & \- & \- & \- Overall\textcolor{bgreen}{\textbf{($\sigma$)}} & Easy &  Hard &   View-dep. &  View-indep.&\- &  Overall\textcolor{bgreen}{\textbf{($\sigma$)}}  &  Easy &  Hard &   View-dep. &  View-indep.\\

    \cmidrule(r){1-2}
    \cmidrule(r){3-4}
    \cmidrule(r){5-9}
    \cmidrule(r){10-15}
    ReferIt3D \citep{achlioptas2020referit3d}       &\- & N/A & \- &  35.6 &  43.6 & 27.9 &  32.5 &  37.1 &\-&  40.8 &  44.7 & 31.5 &  39.2 &  40.8\\
    Text-Guided-GNNs \citep{huang2021text} &\- & N/A & \- &   37.3 &  44.2 & 30.6 &  35.8 &  38.0&\-&  45.0 &  48.5 & 36.9 &  45.8 &  45.0\\
    InstanceRefer  \citep{yuan2021instancerefer}  &\- & N/A & \- &  38.8 &  46.0 & 31.8 &  34.5 &  41.9&\-&  48.0 &  51.1 & 40.5 &  45.4 &  48.1\\
    3DRefTrans \citep{abdelreheem20223dreftransformer} & \- & N/A & \- &  39.0 &  46.4 & 32.0 &  34.7 &  41.2&\-&  47.0 &  50.7 & 38.3 &  44.3 &  47.1\\
    3DVG-Trans \citep{zhao20213dvg} &\- & N/A & \- &  40.8 &  48.5 & 34.8 &  34.8 &  43.7&\-&  51.4 &  54.2 & 44.9 &  44.6 &  51.7\\
    FFL-3DOG   \citep{feng2021free}      & \- & N/A & \- &  41.7 &  48.2 & 35.0 &  37.1 &  44.7&\- & - & - & - & - & -\\  
    TransRefer3D  \citep{he2021transrefer3d}    & \- & N/A & \- & 42.1 &  48.5 & 36.0 &  36.5 &  44.9&\-&  57.4 &  60.5 & 50.2 &  49.9 &  57.7\\
    LanguageRefer \citep{roh2022languagerefer}   & \- & N/A & \- &  43.9 &  51.0 & 36.6 &  41.7 &  45.0 &\- &  56.0 &  58.9 & 49.3 &  49.2 &  56.3\\
    3D-SPS \citep{luo20223d}   & \- & N/A & \- &  51.5 &  58.1 &  45.1 & 48.0 &  53.2 & \-& 62.6 &  56.2 &  65.4 & 49.2 &  63.2 \\
    \hline
    LAR \citep{bakr2022look} &\- & N/A & \- &  48.9 &  56.1 & 41.8 &  46.7 &  50.2 &\- &  59.4     &  63.0  & 51.2  &  50.0     &  59.1\\ 
    \rowcolor{blue!15} {+\papernameAbbrev} (\textbf{Ours})   &\- & \xmark & \- & \textbf{52.5}\textcolor{bgreen}{$\pm$0.1} &  \textbf{59.4} & \textbf{45.7} &  \textbf{50.4} &  \textbf{53.6} &\- &  - &  -     & -     &  -     &  - \\ 
    \rowcolor{blue!15} {+\papernameAbbrev} (\textbf{Ours})   &\- & \cmark & \- &  \textbf{58.0}\textcolor{bgreen}{$\pm$0.2} &  \textbf{65.1} & \textbf{50.9} &  \textbf{54.6} &  \textbf{60.0} &\- &  \textbf{69.0}\textcolor{bgreen}{$\pm$1.0}     &  \textbf{74.2}     & \textbf{59.4}     &  \textbf{58.9}     &  \textbf{70.1}\\ 
    \hline
    
    SAT  \citep{yang2021sat}           & \- & N/A & \- &  49.2 &  56.3 & 42.4 &  46.9 &  50.4\-& &  57.9 &  61.2 & 50.0 &  49.2 &  58.3\\
    +PhraseRefer \citep{yuan2022toward} &\- & \cmark & \- &  54.4 &  62.1 & 47.0 &  51.2 &  56.0 &\- &  63.0     &  66.2  & 55.4  &  48.7  &  63.6 \\
    +ScanEnts \citep{abdelreheem2022scanents3d} &\- & \cmark & \- &  52.5 &  59.8 & 45.6 &  51.3 &  53.2 &\- &  -     &  -  & -  &  -  &  - \\
    \rowcolor{blue!15} {+\papernameAbbrev} (\textbf{Ours})   &\- & \xmark & \- &  \textbf{53.1}\textcolor{bgreen}{$\pm$0.2} &  \textbf{60.3} & \textbf{46.2} &  \textbf{51.8} &  \textbf{53.9} &\- &    - &  -     & -     &  -     &  - \\ 
    \rowcolor{blue!15} {+\papernameAbbrev} (\textbf{Ours})   &\- & \cmark & \- &  \textbf{58.1}\textcolor{bgreen}{$\pm$0.4} &  \textbf{65.1} & \textbf{51.6} &  \textbf{57.0} &  \textbf{58.4} &\- &  \textbf{69.1}\textcolor{bgreen}{$\pm$0.5}     &  \textbf{71.9}     & \textbf{62.8}     &  \textbf{54.5}     &  \textbf{70.2}\\ 
    \hline
    
    MVT \citep{huang2022multi}   & \- & N/A & \- &  55.1 &  61.3 & 49.1 &  54.3 &  55.4 \-& &  64.5 &  66.9 & 58.8 &  58.4 &  64.7\\
    +ScanEnts \citep{abdelreheem2022scanents3d} &\- & \cmark & \- &  59.3 &  65.4 & 53.5 &  57.3 &  60.4 &\- &  -     &  -  & -  &  -  &  - \\
    +PhraseRefer \citep{yuan2022toward} &\- & \cmark & \- &  59.0 &  - & - &  - &  - &\- &  68.0     &  -  & -  &  -  &  - \\
    \rowcolor{blue!15} {+\papernameAbbrev} (\textbf{Ours})   &\- & \xmark & \- &  \textbf{60.4}\textcolor{bgreen}{$\pm$0.2} &  \textbf{66.2} & \textbf{54.6} &  \textbf{58.1} &  \textbf{61.5} &\- &    - &  -     & -     &  -     &  - \\ 
    \rowcolor{blue!15} {+\papernameAbbrev} (\textbf{Ours})   &\- & \cmark & \- &  \textbf{64.4}\textcolor{bgreen}{$\pm$0.2} &  \textbf{70.0} & \textbf{59.2} &  \textbf{61.9} &  \textbf{65.7} &\- &  \textbf{73.2}\textcolor{bgreen}{$\pm$0.5}     &  \textbf{75.2} & \textbf{67.9} &  \textbf{67.6} &  \textbf{73.5}\\ 
    \hline
    
    ViL3DRel \citep{chen2022language}   & \- & N/A & \- &  63.6 &  70.0 & 56.9 &  61.3 &  64.1 \-& &  72.8 &  72.5 & 66.3 &  61.5 &  72.8\\
    \rowcolor{blue!15} {+\papernameAbbrev} (\textbf{Ours})   &\- & \xmark & \- & \textbf{64.1}\textcolor{bgreen}{$\pm$0.1} &  \textbf{70.4} & \textbf{57.5} &  \textbf{61.7} &  \textbf{64.8} &\- &    - &  -     & -     &  -     &  - \\ 
    \rowcolor{blue!15} {+\papernameAbbrev} (\textbf{Ours})   &\- & \cmark & \- &  \textbf{64.0}\textcolor{bgreen}{$\pm$0.2} &  \textbf{70.4} & \textbf{57.3} &  \textbf{61.5} &  \textbf{64.8} &\- &  \textbf{73.6}\textcolor{bgreen}{$\pm$0.5}     &  \textbf{73.0} & \textbf{67.2} &  \textbf{61.9} &  \textbf{73.7}\\ 
  \bottomrule
\end{tabular}
}
\vspace{-3mm}
\label{tab:benchmark_nr3d}
\end{table}
% -------------------------------------------------------

% -------------------------------------------------------
\begin{table}[!tb]
\begin{minipage}[c]{0.63\textwidth}
\captionsetup{font=small}
\centering
\captionof{table}{Benchmarking results while training jointly on Nr3D and Sr3D datasets \citep{achlioptas2020referit3d}.}
\scalebox{0.60}{
\begin{tabular}{ccccccc}
    \toprule
    \multicolumn{1}{c}{\multirow{2}{*}{Method}} & \- & \multicolumn{5}{c}{Nr3D+Sr3D ↑} \\
    \cmidrule(r){3-7}
    \- & \-  &  Overall\textcolor{bgreen}{\textbf{($\sigma$)}} & Easy &  Hard &   View-dep. &  View-indep. \\
    \cmidrule(r){1-2}
    \cmidrule(r){3-7}
    ReferIt3D \citep{achlioptas2020referit3d}      &\-  &  37.2 $\pm$ 0.3 &  44.0 $\pm$ 0.6 & 30.6 $\pm$ 0.3 &  33.3 $\pm$ 0.6 &  39.1 $\pm$ 0.2   \\
    non-SAT  \citep{yang2021sat} &\- &  43.9 $\pm$ 0.3 &  - & - &  - &  -   \\
    TransRefer3D  \citep{he2021transrefer3d}    & \- &  47.2 $\pm$ 0.3 &  55.4 $\pm$ 0.5 & 39.3 $\pm$ 0.5 &  40.3 $\pm$ 0.4 & 50.6 $\pm$ 0.2\\
    SAT  \citep{yang2021sat} & \- &  53.9 $\pm$ 0.2 &  61.5 $\pm$ 0.1 & 46.7 $\pm$ 0.3 &  52.7 $\pm$ 0.7 & 54.5 $\pm$ 0.3\\
    MVT \citep{huang2022multi}  & \-  &  58.5 $\pm$ 0.2 &  65.6 $\pm$ 0.2 & 51.6 $\pm$ 0.3 &  56.6 $\pm$ 0.3 & 59.4 $\pm$ 0.2\\
    \rowcolor{blue!15} {MVT+\papernameAbbrev} (\textbf{Ours})   &\- &  \textbf{62.5}\textcolor{bgreen}{$\pm$0.4} &  \textbf{70.0$\pm$0.5} & \textbf{53.0$\pm$0.2} &  \textbf{58.3$\pm$0.1} &  \textbf{63.9$\pm$0.3}  \\ 
  \bottomrule
\end{tabular}
}
\label{tab_benchmark_joint_nr3d_sr3d}
\end{minipage}
\begin{minipage}[c]{0.37\textwidth}
\captionsetup{font=small}
\centering
\captionof{table}{Benchmarking results on ScanRefer dataset \citep{chen2020scanrefer}.}
\scalebox{0.65}{
\begin{tabular}{cccccc}
    \toprule
    \multicolumn{1}{c}{\multirow{2}{*}{Method}} & \- & \multicolumn{4}{c}{Data Percentage} \\
    \cmidrule(r){3-6}
    \-  & \- & 10\% & 40\% & 70\% & 100\% \\
    \cmidrule(r){1-2}
    \cmidrule(r){3-6}
    MVT  \citep{huang2022multi}     &\- & 36.4 & 51.5 & 54.5 & 57.4 \\
    +PhraseRefer \citep{yuan2022toward} &\- & - & - & - & 59.9 \\
    \rowcolor{blue!15} {+\papernameAbbrev} (\textbf{Ours})   &\- &  \textbf{48.6} &  \textbf{60.1} & \textbf{62.5} &  \textbf{64.2} \\
    \hline
    SAT  \citep{he2021transrefer3d}    & \- & 34.9 & 49.3 & 50.5 & 53.8 \\
    +PhraseRefer \citep{yuan2022toward} &\- & - & - & - & 57.5 \\
    \rowcolor{blue!15} {+\papernameAbbrev} (\textbf{Ours})   &\- &  \textbf{49.1} &  \textbf{57.7} & \textbf{58.9} &  \textbf{60.3}  \\
  \bottomrule
\end{tabular}
}
\label{tab_benchmark_scanrefer}
\end{minipage}
\end{table}
% -------------------------------------------------------

We verify the effectiveness of our proposed framework, termed {\papernameAbbrev} on three well-known 3D visual grounding benchmarks, i.e., Nr3D, Sr3D \citep{achlioptas2020referit3d} and ScanRefer \citep{chen2020scanrefer}.
By effectively localizing a chain of anchors before the final target, we achieve state-of-the-art results without requiring any additional manual annotations.
As shown in Table \ref{tab:benchmark_nr3d}, when we integrate our module into four baselines; LAR \citep{bakr2022look}, SAT \citep{yang2021sat}, MVT \citep{huang2022multi} and ViL \citep{chen2022language}, it boosts the accuracy by 3.6\%, 4\%, 5\%, 0.5\% on Nr3D and 10\%, 11\%, 9\%, 1\% on Sr3D, respectively.
The disparity between the gain achieved on Nr3D and Sr3D is due to our pseudo label module that hinders achieving more gain on Nr3D.
Exchanging of our noisy pseudo labels with precise manual annotations led to a noteworthy enhancement in referring accuracy, where our module boosts the performance by 9\% compared to the baselines; LAR, SAT and MVT, respectively.
This outcome underscores our model's ability to yield enhanced performance not only for simpler descriptions (Sr3D) but also in the context of more intricate, free-form descriptions (Nr3D).
We have a limited gain on only ViL. A detailed analysis that justifies this behaviour is mentioned in the Appendix \ref{appendix_Justification_vil}.

\noindent \textbf{Nr3D+Sr3D.} 
In addition, we jointly train on Nr3D and Sr3D. 
Whereas, we augment the Nr3D training data with Sr3D, while testing on the same original Nr3D test-set.
Consistently with the previous results mentioned in Table \ref{Table_ablation_for_cot_components} and \ref{tab:benchmark_nr3d}, we surpass all the existing work by 3.5\%.
As shown in Table \ref{tab_benchmark_joint_nr3d_sr3d}, we achieve 62.5\% grounding accuracy, while MVT only achieves 58.5\% accuracy.

\noindent \textbf{ScanRefer.} 
To further show the effectiveness of our proposed method, we have conducted several experiments on the ScanRefer dataset across different data percentages on MVT and SAT baselines.
As shown in Table \ref{tab_benchmark_scanrefer}, we outperform both MVT and SAT by a significant margin across the entire data percentages, i.e., 10\%, 40\%, 70\%, and 100\%. 
More specifically, integrating our CoT framework into MVT has boosted the performance by 12.2\%, 8.6\%, 8\%, and 6.8\%, respectively.
In addition to MVT, we have integrated our CoT framework into the SAT baseline, where the same performance gain has been achieved, probing our method's effectiveness across a comprehensive range of baseline models, datasets, and different available data percentages.
% --------------------------------------------------------
% \subsection{Data Efficiency}

\noindent \textbf{Data Efficiency.} 
To further validate the effectiveness of our framework, we assess it on a more challenging setup, where we assume access to only limited data.
Four percentage of data is tested, i.e., 10\%, 40\%, 70\%, and 100\%.
As shown in Figure \ref{fig_data_eff}, on the Sr3D dataset, using only 10\% of the data, we match the same performance of MVT and SAT that are trained on 100\% of the data.
This result highlights the data efficiency of our method.
Furthermore, when trained on 10\% of the data on Nr3D with noisy pseudo labels (Sec. \ref{sec_pseudo_labels}), we still surpass all the baselines with considerable margins.

% \subsection{Qualitative results.}
% --------------------------------------------------------
\begin{figure}[t!]
\captionsetup{font=small}
\centering
\includegraphics[width=0.99\textwidth]{Figures/CoT_Qualitative_results.pdf}
\caption{Our qualitative results. First row indicates the GT w.r.t the input utterance, demonstrated in the last row. Second and third rows show the qualitative results for MVT and our method, respectively. The success and the failure cases are shown in green and red boxes, respectively.}
\vspace{-3mm}
\label{fig_qualitative_results}
\end{figure}
% --------------------------------------------------------
\noindent \textbf{Qualitative results.} As shown in Figure \ref{fig_qualitative_results}, the first three examples show that our model successfully localizes the referred objects by leveraging the mentioned anchors, such as ``the table with 5 chairs around''.
However, in the ambiguous description shown in the fourth example: ``2nd stool from the left'', the model incorrectly predicts the stool, as it is view-dependent. 
In other words, if you look at the stools from the other side, our predicted box will be correct.
Additionally, the last example shows a challenging scenario where a negation in the description is not properly captured by our model.
% --------------------------------------------------------